\begin{table}[tb]
\centering\small\setlength{\tabcolsep}{4pt}
\caption{既知４課題の全８条件．M：実測残時間表示，H：非表示．費用は全論理入出力token．}
\begin{tabular}{llrrr}
課題・条件 & 成否 & 費用 & 生成 & API \\ \hline
写真 M & 成功 & 177,373 & 17 & 100 \\
写真 H & 成功 & 135,628 & 16 & 99 \\
精算 M & 成功 & 284,294 & 20 & 26 \\
精算 H & 成功 & 211,110 & 24 & 25 \\
返金 M & 成功 & 216,854 & 21 & 25 \\
返金 H & 成功 & 102,087 & 13 & 13 \\
書出し M & 未達 & 108,571 & 12 & 80 \\
書出し H & 未達 & 331,569＋不明1 & 20 & 69 \\
\hline
\end{tabular}
\end{table}
Mは4条件中3成功，Hは4条件中3成功だった．全論理入出力はMが787,092，Hが780,394＋不明1である．
未評価または使用量不明を含むため，全費用差は確定しない．
事前の開発上の採用基準は満たさなかった．単一実行の合計差を平均的改善や一般的な品質維持とは解釈しない．
作業時間の合計はMが1,897.7秒，Hが2,975.3秒であり，論理tokenとは別に報告する．
全条件合計は143生成，437 API，17 API例外である．所有推論器の滞在時間は4,935.6秒，この内に含まれる独立評価の時間は13.60秒だった．
表示説明文の差により，Hの初回入力は全課題でMより12 token多い．従って入力文の単純な短縮を評価した比較ではない．各条件の入出力，cache入力，準備・作業・後片付け時間は全件表に保存する．
